\section{Especificación detallada de casos de uso}
\label{sec:casos}

Cada caso de uso sigue el formato extendido de Cockburn (actor, objetivo,
precondiciones, flujo principal, alternativos, postcondiciones) y se enlaza
con los requerimientos que satisface. La cobertura completa se resume en
la matriz de trazabilidad al final de la sección.

\subsection{\cu{01}: Realizar entrevista individual}
\begin{cajaCriterio}
\textbf{Actor primario:} Candidato (\sk{1}, \sk{2}). \\
\textbf{Actores secundarios:} LLM Entrevistador, LLM Coach. \\
\textbf{Objetivo:} Practicar una entrevista de 30 min y obtener un
resumen con plan de mejora. \\
\textbf{Precondiciones:} Usuario autenticado; consentimiento de audio,
video y métricas otorgado; navegador con permisos de cámara y micrófono. \\
\textbf{Postcondiciones (éxito):} Sesión persistida; métricas calculadas;
plan de mejora generado y mostrado al usuario.
\end{cajaCriterio}

\paragraph{Flujo principal}
\begin{enumerate}
  \item El candidato selecciona industria, rol y nivel objetivo.
  \item El sistema valida los permisos del navegador y carga la sala
        virtual.
  \item El avatar entrevistador saluda por voz y formula la primera
        pregunta.
  \item El candidato responde por voz; MediaPipe deriva métricas en el
        navegador y las envía al backend.
  \item El aura luminosa actualiza los indicadores con latencia
        $\le 500\,\mathrm{ms}$.
  \item El entrevistador formula la siguiente pregunta adaptándose a la
        respuesta previa. Se repite hasta completar la sesión o hasta
        que el candidato diga ``terminar entrevista''.
  \item El LLM Coach genera el resumen y plan de mejora; el sistema lo
        muestra al candidato y lo persiste.
\end{enumerate}

\paragraph{Flujos alternativos}
\begin{itemize}
  \item \textbf{4a.} Si MediaPipe no detecta rostro durante $>10\,\mathrm{s}$,
        el sistema pausa y pide reencuadre (sin penalizar postura).
  \item \textbf{6a.} Si el candidato dice ``pausa'', el sistema detiene
        captura y temporizador hasta el comando ``continuar''.
  \item \textbf{6b.} Si el candidato dice ``repite la pregunta'', el
        avatar la reformula sin contarlo como reintento.
  \item \textbf{2a / 7a.} Si falla la API del LLM, el sistema cae en una
        cola de reintentos y avisa al usuario sin perder la sesión.
\end{itemize}

\paragraph{Cubre} \rf{01}--\rf{07}, \rf{15}, \rnf{01}--\rnf{02},
\rnf{05}, \rnf{14}.

\subsection{\cu{02}: Configurar perfil y preferencias de accesibilidad}
\begin{cajaCriterio}
\textbf{Actor primario:} Cualquier usuario primario (\sk{1}--\sk{7}). \\
\textbf{Objetivo:} Establecer industria preferida, idioma, paciencia del
entrevistador, modo de baja estimulación, modo audio-first, contraste,
tipografía, declaración de condiciones que reemplazan métricas. \\
\textbf{Precondiciones:} Usuario autenticado.
\end{cajaCriterio}

\paragraph{Flujo principal}
\begin{enumerate}
  \item El usuario abre ``Mi perfil''.
  \item Selecciona o ajusta cada preferencia; la UI muestra una vista
        previa en vivo del impacto.
  \item Guarda; las preferencias se aplican inmediatamente a sesiones
        futuras y al asistente lateral.
\end{enumerate}

\paragraph{Cubre} \rf{17}--\rf{20}, \rnf{08}--\rnf{09}.

\subsection{\cu{03}: Realizar mock interview entre pares (peer)}
\begin{cajaCriterio}
\textbf{Actores primarios:} Candidato (\sk{1}) y Peer Observador (\sk{3}). \\
\textbf{Objetivo:} Practicar una entrevista con feedback humano sin
salir de la aplicación. \\
\textbf{Precondiciones:} Ambos usuarios autenticados; uno de ellos creó
una sala con código compartible; ambos consintieron grabación.
\end{cajaCriterio}

\paragraph{Flujo principal}
\begin{enumerate}
  \item El candidato crea sala; el sistema emite un código.
  \item El peer entra con el código; WebRTC establece la conexión.
  \item El candidato responde a preguntas que pueden venir del LLM
        entrevistador, del peer, o de ambos según el modo.
  \item El peer ancla comentarios en la línea de tiempo durante o
        después de cada respuesta.
  \item Al terminar, los usuarios pueden intercambiar roles.
  \item El sistema genera dos resúmenes (uno por rol) y los persiste.
\end{enumerate}

\paragraph{Flujos alternativos}
\begin{itemize}
  \item \textbf{2a.} Si la conexión WebRTC falla, el sistema reintenta
        con servidor TURN; si persiste, ofrece modo asincrónico
        (grabar y compartir).
\end{itemize}

\paragraph{Cubre} \rf{10}--\rf{12}, \rf{21}, \rnf{04}, \rnf{15}.

\subsection{\cu{04}: Consultar historial y plan de mejora}
\begin{cajaCriterio}
\textbf{Actor primario:} Candidato (\sk{1}, \sk{2}). \\
\textbf{Actor secundario:} Coach o mentor (\sk{9}) si el usuario lo
autorizó. \\
\textbf{Objetivo:} Revisar la evolución por competencia, las sesiones
pasadas y el plan de mejora vigente.
\end{cajaCriterio}

\paragraph{Flujo principal}
\begin{enumerate}
  \item El usuario abre ``Mi progreso''.
  \item Visualiza tendencias por competencia (gráfico) y \textit{badges}
        obtenidos.
  \item Selecciona una sesión específica para revisar la grabación,
        las métricas y los comentarios.
  \item Pide al asistente lateral que profundice en un punto del plan
        (``dame ejemplos para reducir muletillas''); el LLM responde con
        contexto del historial.
\end{enumerate}

\paragraph{Cubre} \rf{07}--\rf{09}, \rf{16}, \rnf{14}.

\subsection{\cu{05}: Gestionar gamificación (rangos, badges, ligas)}
\begin{cajaCriterio}
\textbf{Actor primario:} Candidato (\sk{1}--\sk{3}). \\
\textbf{Objetivo:} Inscribirse en una liga, ver ranking, recibir
\textit{badges} al cumplir hitos.
\end{cajaCriterio}

\paragraph{Flujo principal}
\begin{enumerate}
  \item El usuario explora ligas activas filtrando por industria.
  \item Se inscribe en la liga correspondiente.
  \item Cada sesión válida actualiza su puntaje semanal.
  \item Al cumplir el hito de un \textit{badge} (ej. 3 sesiones con
        cero muletillas en 5 min), el sistema lo otorga y notifica.
  \item Al cierre semanal, se publica el ranking y el usuario puede
        ver su posición.
\end{enumerate}

\paragraph{Cubre} \rf{13}--\rf{14}, \rnf{04}.

\subsection{\cu{06}: Gestionar consentimiento, datos y privacidad}
\begin{cajaCriterio}
\textbf{Actor primario:} Cualquier usuario primario (\sk{1}--\sk{7}). \\
\textbf{Actor regulador:} Autoridad de protección de datos (\sk{16}). \\
\textbf{Objetivo:} Otorgar/revocar consentimientos, eliminar grabaciones,
exportar datos.
\end{cajaCriterio}

\paragraph{Flujo principal}
\begin{enumerate}
  \item En el primer ingreso, el sistema solicita consentimiento granular
        (audio, video, transcripción, métricas, uso comunitario).
  \item En cualquier momento, el usuario puede revocar un consentimiento
        específico desde ``Mi privacidad''.
  \item El usuario solicita borrado de una grabación; el sistema lo
        ejecuta y emite confirmación.
  \item El usuario solicita exportación de datos; el sistema genera un
        archivo cifrado y le envía un enlace temporal.
\end{enumerate}

\paragraph{Cubre} \rf{21}--\rf{22}, \rnf{05}--\rnf{07}.

\subsection{\cu{07}: Administrar la plataforma}
\begin{cajaCriterio}
\textbf{Actor primario:} Administrador del sistema (\sk{12}). \\
\textbf{Objetivo:} Gestionar usuarios, contenido de industrias, ligas y
moderación del feed comunitario.
\end{cajaCriterio}

\paragraph{Flujo principal}
\begin{enumerate}
  \item El administrador entra al \textit{backoffice}.
  \item Crea, edita o desactiva ligas, industrias o usuarios.
  \item Revisa el feed comunitario; modera o elimina contenido reportado.
\end{enumerate}

\paragraph{Cubre} \rf{23}.

\subsection{\cu{08}: Generar reportes institucionales}
\begin{cajaCriterio}
\textbf{Actor primario:} Universidad / oficina de empleabilidad (\sk{10}). \\
\textbf{Objetivo:} Obtener métricas agregadas anonimizadas para reportes
de empleabilidad.
\end{cajaCriterio}

\paragraph{Flujo principal}
\begin{enumerate}
  \item Un usuario institucional autenticado accede al panel de reportes.
  \item Filtra por industria, rango de fechas y rangos por competencia.
  \item Descarga un reporte en CSV/PDF con datos agregados (nunca
        identificables).
\end{enumerate}

\paragraph{Cubre} \rf{24}, \rnf{05}.

\subsection{Matriz de trazabilidad caso de uso a requerimiento}

\begin{longtable}{p{1.6cm} p{12cm}}
\toprule
\textbf{CU} & \textbf{Requerimientos cubiertos} \\
\midrule
\cu{01} & \rf{01}--\rf{07}, \rf{15}; \rnf{01}, \rnf{02}, \rnf{05}, \rnf{14} \\
\cu{02} & \rf{17}--\rf{20}; \rnf{08}, \rnf{09} \\
\cu{03} & \rf{10}--\rf{12}, \rf{21}; \rnf{04}, \rnf{15} \\
\cu{04} & \rf{07}--\rf{09}, \rf{16}; \rnf{14} \\
\cu{05} & \rf{13}--\rf{14}; \rnf{04} \\
\cu{06} & \rf{21}--\rf{22}; \rnf{05}--\rnf{07} \\
\cu{07} & \rf{23} \\
\cu{08} & \rf{24}; \rnf{05} \\
\bottomrule
\end{longtable}

\noindent Los RNF transversales \rnf{03}, \rnf{08}, \rnf{10}--\rnf{13} se
verifican a nivel de toda la aplicación (rendimiento global, accesibilidad
en todas las vistas, mantenibilidad, observabilidad y costo) y no se
asignan a un caso de uso individual.
